\documentclass{article}
\usepackage{graphicx} % Required for inserting images
\usepackage[utf8]{inputenc}
\usepackage{amsmath}

\title{CTA200H Assignment 3}
\author{Aniketh Salil}
\date{May 2026}

\begin{document}

\maketitle
\section{Introduction}

\section{Methods}
\subsection{Complex Plane Iteration; Mandelbrot Set}
The Mandelbrot Set is an famous set of numbers in the complex plane that produce a "fractal" (a shape with infinite perimeter and finite area). Essentially, in a complex plane where the boundaries are
\begin{equation}
    -2 < x < 2 \text{ and} -2 < y < 2
\end{equation}
there are a set of points that do not diverge after a set number of iterations of the following formulae:

\begin{equation}
    c = x + iy
\end{equation}
\begin{equation}
    z_{i+1} = z_i + c
\end{equation}
Where $z_n$ is the $n^{th}$ iteration of equation 2.

Using numpy, we can perform the iteration for points within the boundary and observe whether or not they diverge within a fixed number of iterations. We then plotted the points all points within the boundary by using the matplotlib.pyplot function, "imshow", which displays points in a complex plane. From these points, two plots were generated; one that showed which points diverged, and another that showed how many iterations it took for divergence to occur (or if the maximum number of iterations was reached).

\subsection{Deterministic Modeling of Earth's Atmosphere}
From Lorenz's paper, we recreated his modelling of the Earth's atmosphere considering thermal convection and fluid viscosity. The equations to model this are:
\begin{equation}
\dot{X} = -\sigma(X - Y)
\end{equation}
\begin{equation}
\dot{Y} = rX - Y - XZ
\end{equation}
\begin{equation}
\dot{Z} = -bZ + XY
\end{equation}
Using the solve\_ivp function from the scipy.integrate library, we can find a solution to this ODE, and observe the behaviour of each parameter (X, Y, and Z) over a number of iterations.
Furthermore, we can use numpy to get slighly altered initial conditions, where W' = (0., 1.00000001, 0.), instead of W = (0, 1, 0). Then, we find the distance between W' and W over time. 
\section{Results}
\subsection{Mandelbrot Plots}

\begin{figure}[h]
    \centering
    \includegraphics[width=1.2\linewidth]{mandelbrot.png}
    \caption{Mandelbrot Plots with boundaries; $-2 < x < 2$ and $-2 < y< 2$}
    \label{fig:placeholder}
\end{figure}
Here, we can observe that there is a coherent "fractal" shape that is formed along the boundary of points that diverge (purple), and do not diverge (yellow) after a number of iterations. In the second plot, we can observe a few distinct features. The black represents points that diverge above |2| within a small number of iterations (<<10), whereas the "warmer" colours represent points that diverge in a greater number of iterations. Finally, the yellow points within the fractal are points that do not diverge after 100 iterations.

\subsection{Lorenz Plots}
\begin{figure}
    \centering
    \includegraphics[width=1\linewidth]{Y plots Lorenz.png}
    \caption{Recreated Figure 1 of Lorenz's plots. The x-axis represents the number of iterations, and the y-axis is the corresponding y-value}
    \label{fig:placeholder}
\end{figure}
In this figure, we can observe various epochs within the plot. Initially, there is a dramatic spike, before it oscillates. After roughly 1600 iterations, there seem to be numerous, ordered, but significant spikes 

\begin{figure}
    \centering
    \includegraphics[width=1\linewidth]{xy and yz.png}
    \caption{X-Y and Y-Z plots}
    \label{fig:placeholder}
\end{figure}
We can also observe cyclic behaviour over time in the X-Y and Y-Z plots above. 
\begin{figure}
    \centering
    \includegraphics[width=1\linewidth]{W'plots.png}
    \caption{W' distance plots}
    \label{fig:placeholder}
\end{figure}
In the above plot, we measured the distance between W' and W. This behaviour is consistent with our predictions as the plot follows a straight line initially, but starts to oscillate with growing amplitude over time. This is likely because of the cyclic nature of the system observed in the previous plots. With slight variations in the initial conditions, it is expected that they will slowly diverge over time.
\end{document}
